% GENERATED by assurance/reports/build_reports.py from the records; do not edit.
\documentclass[11pt,a4paper]{report}
\usepackage{../latex/liondriver-report}
% GENERATED by assurance/reports/build_reports.py from the records; do not edit.
\newcommand{\FscVersion}{0.32}
\newcommand{\FscDate}{2026-11-01}
\newcommand{\GoalCount}{8}
\newcommand{\FsrCount}{58}
\newcommand{\FsrAsilD}{0}
\newcommand{\FsrAsilC}{30}
\newcommand{\FsrAsilB}{13}
\newcommand{\FsrAsilA}{6}
\newcommand{\FsrQM}{4}
\newcommand{\FsrExt}{4}
\newcommand{\FsrRef}{1}

\LDsetdoc{WP-C-04 Functional Safety Concept}{\FscVersion}
\title{Functional Safety Concept\\[4pt]\large WP-C-04, version \FscVersion}
\author{LionDriver assurance\\\small generated from \texttt{assurance/02-concept/WP-C-04-functional-safety-concept.md}}
\date{\FscDate}
\begin{document}
\maketitle
\tableofcontents
\chapter{Introduction}
This report renders the functional safety concept of LionDriver (work product WP-C-04, version \FscVersion{} of \FscDate{}) from its source of record, as ISO 26262-3 asks of the concept phase \cite{iso26262-3}. It derives \FsrCount{} functional safety requirements (FSRs) from \GoalCount{} safety goals: \FsrAsilC{} at ASIL~C, \FsrAsilB{} at ASIL~B, \FsrAsilA{} at ASIL~A and \FsrQM{} at QM, with \FsrExt{} external measures outside the item and \FsrRef{} cross-reference to other FSRs.

Every number in this report is a macro written from the record (\texttt{numbers.tex}), and every table is read from it, so the report cannot fall out of step with the analysis: when the record changes, the build fails until the report is regenerated. The record is a draft, AI-assisted (independence I0) and not reviewed; so is this report.
\chapter{Safety goals}
Table~\ref{tab:goals} summarises the safety goals of the hazard analysis (WP-C-03) that the concept refines. A \warn{} marks an ASIL that rises if an assumption of use fails.
{\small
\begin{longtable}{@{}p{1.2cm}p{4.3cm}p{2.3cm}p{3.6cm}p{2.9cm}@{}}
\caption{Safety goals with their ASIL, safe state and preliminary FTTI (WP-C-04 \S3).}\label{tab:goals}\\
\toprule \textbf{SG} & \textbf{Statement} & \textbf{ASIL} & \textbf{Safe state} & \textbf{FTTI} \\ \midrule\endfirsthead
\toprule \textbf{SG} & \textbf{Statement} & \textbf{ASIL} & \textbf{Safe state} & \textbf{FTTI} \\ \midrule\endhead
SG-01 & No uncontrollable lateral motion, no steering actuation while not engaged & C \warn{} (D if AOU-01/02 fail) & LKA torque 0, steer request cleared & $\le$ 0.5 s \\ \midrule
SG-02 & No loss/degradation of lateral control without adequate take-over warning & C & Warning + torque ramp-down & $\le$ 1 s to warning \\ \midrule
SG-03 & No unintended acceleration & C \warn{} (D if AOU-03/05 fail) & Accel $\le$ 0; ACC cancelled & $\le$ 1 s \\ \midrule
SG-04 & No deceleration beyond what following traffic and driver can control & B \warn{} (C if AOU-05 fails) & Inactive (no ACC command) with bounded jerk & $\le$ 1 s \\ \midrule
SG-05 & Release on driver brake, cancel, steering override & B \warn{} (C if AOU-01/02/03 fail) & Disengaged & $\le$ 0.2 s \\ \midrule
SG-06 & No loss of longitudinal deceleration without take-over warning & C & Warning; ACC cancelled & $\le$ 1 s \\ \midrule
SG-07 & No suppression, delay or alteration of stock PCS/AEB & B & Camera PCS messages forwarded; relay released if forwarding cannot be guaranteed & $\le$ 0.2 s for forwarding loss (preliminary, WP-S-04 §3.6; accepted as the working budget until measured, DQ-062); content: per frame \\ \midrule
SG-08 & No continued automated control while the driver is not supervising or not responding, without escalation to a minimal-risk reaction & A & Escalating alerts; on no response, controlled deceleration in lane to a standstill with steering on, then stay stopped and held with re-engagement locked out until the driver responds (DQ-365, A) & $\le$ 6 s from loss of attention to the first alert (WP-C-08 OI-2); time to the no-response reaction set in WP-S-04 \\ \midrule
\bottomrule
\end{longtable}
}
\chapter{Functional safety requirements}
Figure~\ref{fig:alloc} shows every safety goal with the FSRs that refine it; the sections below list them goal by goal, with their ASIL, the element they are allocated to and how they are verified. The implementation status of each FSR is kept in the record (WP-C-04 \S5).
\begin{landscape}\begin{figure}[p]\centering
\includegraphics[width=\linewidth,height=0.88\textheight,keepaspectratio]{fig-c04-5-1-sg-fsr-allocation.pdf}
\caption{Safety goals and their FSRs (WP-C-04 Figure 5-1).}\label{fig:alloc}
\end{figure}\end{landscape}
\section{SG-01 Lateral motion (ASIL C \warn{})}
{\small
\begin{longtable}{@{}p{1.7cm}p{7.4cm}p{1.1cm}p{1.8cm}p{2.3cm}@{}}
\caption{FSRs of SG-01 (WP-C-04 \S5.1).}\label{tab:fsr-5.1}\\
\toprule \textbf{ID} & \textbf{Requirement} & \textbf{ASIL} & \textbf{Alloc.} & \textbf{Verification} \\ \midrule\endfirsthead
\toprule \textbf{ID} & \textbf{Requirement} & \textbf{ASIL} & \textbf{Alloc.} & \textbf{Verification} \\ \midrule\endhead
FSR-01.01 & The envelope shall limit the commanded LKA torque magnitude to a speed-dependent bound derived from driver controllability (lateral acceleration and lateral deviation reached within the FTTI). & C & E-03 & A, T-SIL, T-VEH \\ \midrule
FSR-01.02 & The envelope shall limit the rate of change of the commanded LKA torque (rise per frame and over a real-time window) to bounds derived from controllability. A fall toward zero is not limited, as it is the safe direction for SG-01; a bounded ramp of torque withdrawn on a fault is FSR-02.04. & C & E-03 & A, T-SIL \\ \midrule
FSR-01.03 & The envelope shall limit the commanded LKA torque to within a bounded difference of the measured EPS motor torque. & C & E-03 & T-SIL \\ \midrule
FSR-01.04 & The envelope shall block any non-zero LKA torque or set steer-request bit while lateral authority is not granted. & C & E-03 & T-SIL, T-HIL \\ \midrule
FSR-01.05 & The envelope shall grant actuation authority only on a driver engagement action confirmed by the vehicle (rising edge of PCM cruise active) and shall revoke it when PCM cruise becomes inactive. & C & E-03 & T-SIL, T-VEH \\ \midrule
FSR-01.06 & The envelope shall detect loss, corruption, repetition and staleness of each RX signal it uses for gating or limiting (PCM cruise state, brake, gas, driver/EPS torque, wheel speed) and revoke actuation authority within the FTTI allocation of §8. & C & E-03 & T-SIL, FI \\ \midrule
FSR-01.07 & The envelope shall detect missing or stale control commands from the SoC and remove actuation within the FTTI allocation of §8. & C & E-03 & T-HIL, FI \\ \midrule
FSR-01.08 & The envelope shall detect corruption, loss and reordering of command frames received from the SoC and discard affected frames. & C & E-03 & T-HIL, FI \\ \midrule
FSR-01.09 & The safety MCU shall detect its own execution faults (program-flow hang, clock failure, memory corruption, register corruption) and enter the safe state (no actuation TX, relay released) within the FTTI allocation of §8. & C & E-03 (HW+SW) & A (FMEDA), FI \\ \midrule
FSR-01.10 & The safety mode and safety parameter for the reference configuration shall be fixed in the envelope and shall not be changeable by the SoC while a car safety mode is active; debug-only commands shall be unavailable in release builds; and the SoC shall not be able to make the safety MCU run non-release code or drive the relay outside the envelope through its reset and boot lines (GAP-38, refined by TSR-511). & C & E-03 & R, T-HIL \\ \midrule
FSR-01.11 & The envelope shall detect a harness relay malfunction (camera control messages seen on the car side) and inhibit all actuation TX. & C & E-03, E-05 & T-SIL, T-HIL \\ \midrule
FSR-01.12 & On any detected fault affecting lateral actuation, the envelope shall enter the lateral safe state (torque 0, steer request cleared) and keep it until a new valid engagement. & C & E-03 & T-SIL, FI \\ \midrule
FSR-01.13 & The EPS shall limit LKA torque to an overpowerable level and shall end LKA torque within a bounded time after \texttt{0x2E4} stops or carries the request bit cleared. (External measure, AOU-01R) & — (ext.) & EXT-EPS & T-VEH \\ \midrule
FSR-01.14 & The SoC controller shall command lateral motion with a margin below every envelope limit (target: commanded lateral acceleration $\le$ 3.0 m/s$^2$, jerk $\le$ 5 m/s$^3$, torque $\le$ 90 \% of the envelope bound). & QM & E-01 & T-SIL \\ \midrule
\bottomrule
\end{longtable}
}
\section{SG-02 Lateral loss without warning (ASIL C)}
{\small
\begin{longtable}{@{}p{1.7cm}p{7.4cm}p{1.1cm}p{1.8cm}p{2.3cm}@{}}
\caption{FSRs of SG-02 (WP-C-04 \S5.2).}\label{tab:fsr-5.2}\\
\toprule \textbf{ID} & \textbf{Requirement} & \textbf{ASIL} & \textbf{Alloc.} & \textbf{Verification} \\ \midrule\endfirsthead
\toprule \textbf{ID} & \textbf{Requirement} & \textbf{ASIL} & \textbf{Alloc.} & \textbf{Verification} \\ \midrule\endhead
FSR-02.01 & The item shall detect EPS conditions that remove or reduce LKA assistance (EPS LKA fault states) while engaged. & C & E-03 (detection), E-01 (diagnosis text) & T-SIL, T-VEH \\ \midrule
FSR-02.02 & The item shall detect loss of the inputs that the lateral command depends on (driving model output, calibration, localization, vehicle state) and shall not continue to actuate on inputs that are no longer valid. & C & E-01 (QM, see note), E-03 (FSR-01.07 as backstop) & T-SIL, FI \\ \midrule
FSR-02.03 & On a detected loss or degradation of lateral control, the item shall give a visual and acoustic take-over request within 1 s of fault detection and before or at the start of torque reduction. & C & E-01/HMI, E-03 siren for SoC loss & T-HIL, T-VEH \\ \midrule
FSR-02.04 & When torque is withdrawn because of a fault (not a driver action), the reduction shall follow a bounded ramp unless the fault makes the command itself untrusted, in which case torque shall go to zero at once. & C & E-01, E-03 & A, T-SIL \\ \midrule
FSR-02.05 & When granted actuation authority ends for a reason other than a driver action that the safety MCU itself observes, the safety MCU shall give an acoustic warning that does not depend on the SoC. This covers the loss of communication with the SoC or of its commands, a fault the safety MCU detects, a disengagement reported by the SoC whose cruise cancel had no effect, and a reset of the safety MCU (watchdog or fault reset) that interrupts granted authority. The reset case was added 2026-10-14 (DQ-054: the engaged state is kept across the reset; TSR-516(d)); the trigger was widened 2026-10-21 to match TSR-516(b)/(c) (DQ-404 A; cancel-lag guard, lead 2026-10-21). & C & E-03 (panda buzzer) & T-HIL \\ \midrule
FSR-02.06 & Driver monitoring shall detect an inattentive or unresponsive driver while engaged and escalate warnings within the timing set in WP-C-08 §5. & A & E-01/DM, E-02 & T-SIL, T-VEH \\ \midrule
FSR-02.07 & Driver monitoring shall detect its own invalidity (no driver image, camera blocked, model output missing or not credible) and fall back to a policy that is no less strict than the active policy. & A & E-01/DM & T-SIL, FI \\ \midrule
FSR-02.08 & Driver monitoring shall not be disabled, put into demo mode or replaced by synthetic inputs while the item can engage. & A & E-01/DM & R, T-SIL \\ \midrule
FSR-02.09 & If the driver does not respond to the highest DM warning, the item shall decelerate the vehicle in a controlled way in its lane to a standstill, keep it stopped and held, and prevent the item from driving off again or being re-engaged until the driver responds; after the driver responds, the lockout period applies. This is the SG-08 safe state (WP-C-03 v0.4), decided 2026-10-16, DQ-365 A: steering stays on while the vehicle decelerates in lane to a standstill within the envelope limits with the alert sounding; the vehicle then stays stopped and held, and the item cannot drive off again: re-engagement and resumption are locked out until the driver responds. Lateral and longitudinal control shall not stay active at the set speed. & A & E-01/DM & T-SIL, T-VEH \\ \midrule
FSR-02.10 & The item shall detect that the envelope does not grant actuation authority while the SoC is engaged (envelope blocking, stuck authority, safety-mode mismatch) and give a take-over request within the SG-02/SG-06 warning budget (§8). & C & E-01 (detection, alert), E-03 (reports authority state) & T-SIL, FI \\ \midrule
FSR-02.11 & The item shall detect failure of the HMI path that shows the engagement state and gives take-over requests (display, sound) and treat it as a fault leading to disengagement with the SoC-independent indications (FSR-02.05 buzzer, FSR-05.05 cluster cruise state). & C & E-01/HMI, E-03 & T-SIL, FI \\ \midrule
FSR-02.12 & Before granting engagement, and while the item can engage, the item shall read the stored diagnostic trouble codes of the EPS and of the brake and stability control, and shall not allow engagement while a code that affects lateral or longitudinal control is set; the driver shall be told why. (Decided 2026-10-14, DQ-056: replaces reliance on the driver's dashboard lamp check, WP-O-02 OPS-01.) & QM (supports FSR-02.01, FSR-06.03) & E-01 & T-SIL, T-VEH \\ \midrule
\bottomrule
\end{longtable}
}
\section{SG-03 Unintended acceleration (ASIL C \warn{})}
{\small
\begin{longtable}{@{}p{1.7cm}p{7.4cm}p{1.1cm}p{1.8cm}p{2.3cm}@{}}
\caption{FSRs of SG-03 (WP-C-04 \S5.3).}\label{tab:fsr-5.3}\\
\toprule \textbf{ID} & \textbf{Requirement} & \textbf{ASIL} & \textbf{Alloc.} & \textbf{Verification} \\ \midrule\endfirsthead
\toprule \textbf{ID} & \textbf{Requirement} & \textbf{ASIL} & \textbf{Alloc.} & \textbf{Verification} \\ \midrule\endhead
FSR-03.01 & The envelope shall limit the commanded acceleration to an upper bound derived from controllability for the reference vehicle. & C & E-03 & A, T-SIL, T-VEH \\ \midrule
FSR-03.02 & The envelope shall allow only the inactive acceleration value while longitudinal authority is not granted or while the driver presses the accelerator. & C & E-03 & T-SIL \\ \midrule
FSR-03.03 & The envelope shall limit the rate of increase of commanded acceleration (jerk). & C & E-03 & T-SIL \\ \midrule
FSR-03.04 & FSR-01.05\ldots{}01.12 apply to the longitudinal path (shared mechanisms; see trace §9). & C & E-03 & as referenced \\ \midrule
FSR-03.05 & On a detected fault affecting longitudinal control, the item shall stop requesting positive acceleration and request ACC cancel from the PCM. & C & E-03 (inactive value), E-01 (cancel) & T-SIL, T-VEH \\ \midrule
FSR-03.06 & The PCM shall bound ACC acceleration requests to its own envelope and honour the cancel request. (External measure, AOU-05R) & — (ext.) & EXT-PCM & T-VEH \\ \midrule
\bottomrule
\end{longtable}
}
\section{SG-04 Excessive deceleration (ASIL B \warn{})}
{\small
\begin{longtable}{@{}p{1.7cm}p{7.4cm}p{1.1cm}p{1.8cm}p{2.3cm}@{}}
\caption{FSRs of SG-04 (WP-C-04 \S5.4).}\label{tab:fsr-5.4}\\
\toprule \textbf{ID} & \textbf{Requirement} & \textbf{ASIL} & \textbf{Alloc.} & \textbf{Verification} \\ \midrule\endfirsthead
\toprule \textbf{ID} & \textbf{Requirement} & \textbf{ASIL} & \textbf{Alloc.} & \textbf{Verification} \\ \midrule\endhead
FSR-04.01 & The envelope shall limit the commanded deceleration to a bound derived from following-traffic controllability. & B & E-03 & A, T-SIL \\ \midrule
FSR-04.02 & The envelope shall limit the onset rate of commanded deceleration (negative jerk). & B & E-03 & T-SIL \\ \midrule
FSR-04.03 & The transition to the longitudinal safe state shall not create a deceleration step larger than the jerk bound of FSR-04.02. & B & E-01, E-03 & A, T-VEH \\ \midrule
FSR-04.04 & Driver accelerator input shall override commanded deceleration. & B & E-03, EXT-PCM & T-SIL, T-VEH \\ \midrule
FSR-04.05 & Deceleration that the item initiates for a reason other than following traffic shall be announced by a visual and acoustic alert, ahead of its onset for the driver-monitoring no-response and at the latest at its onset for a fault soft disable (DQ-407 A: that slow-down is mild, and holding it back keeps the fault active), and shall stay within the bounds of FSR-04.01/04.02. & B & E-01 (alert), E-03 (bounds) & T-SIL, T-VEH \\ \midrule
FSR-04.06 & Forward collision warnings shall be raised only from valid, fresh inputs, and the false-warning rate in the ODD shall be bounded so that driver reflex braking with a close follower is not a credible consequence. & A & E-01 & T-SIL, validation (WP-V-02) \\ \midrule
\bottomrule
\end{longtable}
}
\section{SG-05 Release on driver action (ASIL B \warn{})}
{\small
\begin{longtable}{@{}p{1.7cm}p{7.4cm}p{1.1cm}p{1.8cm}p{2.3cm}@{}}
\caption{FSRs of SG-05 (WP-C-04 \S5.5).}\label{tab:fsr-5.5}\\
\toprule \textbf{ID} & \textbf{Requirement} & \textbf{ASIL} & \textbf{Alloc.} & \textbf{Verification} \\ \midrule\endfirsthead
\toprule \textbf{ID} & \textbf{Requirement} & \textbf{ASIL} & \textbf{Alloc.} & \textbf{Verification} \\ \midrule\endhead
FSR-05.01 & The envelope shall revoke all actuation authority on a driver brake press (rising edge, or pressed while moving) within the SG-05 FTTI. & B & E-03 & T-SIL, T-VEH \\ \midrule
FSR-05.02 & The envelope shall revoke all actuation authority when the PCM reports cruise inactive (cancel, main switch off). & B & E-03 & T-SIL, T-VEH \\ \midrule
FSR-05.03 & The envelope shall detect driver steering override (driver torque above a threshold) and reduce the LKA torque toward zero at a bounded rate so it does not oppose the driver, for as long as the override lasts; the item shall show the driver that the override was noticed. (Decided 2026-10-14, DQ-053: back off on override and show a cue.) & B & E-03, E-01/HMI (cue) & T-SIL, T-VEH \\ \midrule
FSR-05.04 & The envelope shall block positive acceleration while the accelerator is pressed. & B & E-03 & T-SIL \\ \midrule
FSR-05.05 & The engagement state that the envelope enforces shall equal the PCM cruise state, so that the vehicle's own cruise indicator gives the driver an indication of engagement independent of the item's HMI. & C & E-03, EXT-CLU & A, T-VEH \\ \midrule
FSR-05.06 & If PCM cruise is active while the item is not engaged (or the reverse) for longer than a bounded time, the item shall warn the driver and request cancel. & C & E-01, E-03 & T-SIL \\ \midrule
FSR-05.07 & Driver brake input shall produce braking regardless of ACC commands. (External measure, AOU-03R) & — (ext.) & EXT-PCM & T-VEH \\ \midrule
\bottomrule
\end{longtable}
}
\section{SG-06 Longitudinal loss without warning (ASIL C)}
{\small
\begin{longtable}{@{}p{1.7cm}p{7.4cm}p{1.1cm}p{1.8cm}p{2.3cm}@{}}
\caption{FSRs of SG-06 (WP-C-04 \S5.6).}\label{tab:fsr-5.6}\\
\toprule \textbf{ID} & \textbf{Requirement} & \textbf{ASIL} & \textbf{Alloc.} & \textbf{Verification} \\ \midrule\endfirsthead
\toprule \textbf{ID} & \textbf{Requirement} & \textbf{ASIL} & \textbf{Alloc.} & \textbf{Verification} \\ \midrule\endhead
FSR-06.01 & The item shall detect loss or staleness of the longitudinal plan, lead information (radar/vision) and model output while engaged. & C & E-01 (QM), E-03 backstop via FSR-01.07 & T-SIL, FI \\ \midrule
FSR-06.02 & On detected loss, the item shall give a take-over request within 1 s and cancel ACC so that control returns to the driver; it shall not hold a stale acceleration command. & C & E-01, E-03 & T-SIL, T-VEH \\ \midrule
FSR-06.03 & The item shall detect that the PCM does not execute ACC (ACC fault, cruise fault) and warn the driver. & C & E-01 & T-SIL \\ \midrule
FSR-06.04 & Non-finite or out-of-range actuator commands shall be treated as faults with a fault reaction, not silently replaced. & C & E-01, E-03 & T-SIL \\ \midrule
FSR-06.05 & FSR-02.05 (independent acoustic warning on SoC loss) and FSR-02.06\ldots{}02.09 (DM) also serve SG-06. & as ref. & as ref. & as ref. \\ \midrule
\bottomrule
\end{longtable}
}
\section{SG-07 Stock PCS preservation (ASIL B)}
{\small
\begin{longtable}{@{}p{1.7cm}p{7.4cm}p{1.1cm}p{1.8cm}p{2.3cm}@{}}
\caption{FSRs of SG-07 (WP-C-04 \S5.7).}\label{tab:fsr-5.7}\\
\toprule \textbf{ID} & \textbf{Requirement} & \textbf{ASIL} & \textbf{Alloc.} & \textbf{Verification} \\ \midrule\endfirsthead
\toprule \textbf{ID} & \textbf{Requirement} & \textbf{ASIL} & \textbf{Alloc.} & \textbf{Verification} \\ \midrule\endhead
FSR-07.01 & The item shall forward all camera-side messages to the car side unchanged, except the defined set of intercepted control messages. & B & E-03 & T-SIL, T-HIL \\ \midrule
FSR-07.02 & The item shall not originate any message that the stock PCS function uses (e.g. PRE\_COLLISION \texttt{0x344}, PCS\_HUD \texttt{0x411}, AEB \texttt{0x283} with content) in the reference configuration. & B & E-03 & R, T-SIL \\ \midrule
FSR-07.03 & On safety-MCU fault, loss of device power, or harness fault, the relay shall return to the stock camera-to-vehicle connection. & B & E-05, E-03 & A, T-HIL \\ \midrule
FSR-07.04 & openpilot longitudinal commands shall not prevent PCS braking from taking priority in the PCM. (External measure, AOU-04R) & — (ext.) & EXT-PCM & T-VEH \\ \midrule
FSR-07.05 & The item shall inform the driver when the stock system reports an AEB or FCW event. & QM & E-01 & T-SIL \\ \midrule
FSR-07.06 & If the harness cannot return to the stock camera-to-vehicle connection (switch or relay stuck in intercept, detected by state readback), the item shall inhibit all actuation, warn the driver that stock PCS/AEB is unavailable, and refuse engagement until the harness is repaired. (Accepted 2026-10-14 as the degraded state, DQ-064; applies once readback exists.) & B & E-03 (detection), E-05 (readback), E-01/HMI (warning) & T-HIL, FI \\ \midrule
FSR-07.07 & While the stock system reports an AEB event, the item shall request zero acceleration from the PCM (no added drive), so that stock AEB braking does not depend on PCM arbitration alone. FSR-07.04 and its AOU-04R test still apply. (Decided 2026-10-18, DQ-300 A; System FMEA SFM-036, SFM-ACT-020; FTA MCS-07.10.) & QM (no ASIL claim; FSR-07.04 stays the measure) & E-01 & T-SIL, T-VEH \\ \midrule
\bottomrule
\end{longtable}
}
\section{SG-08 Driver not supervising or not responding (ASIL A)}
{\small
\begin{longtable}{@{}p{1.7cm}p{7.4cm}p{1.1cm}p{1.8cm}p{2.3cm}@{}}
\caption{FSRs of SG-08 (WP-C-04 \S5.8).}\label{tab:fsr-5.8}\\
\toprule \textbf{ID} & \textbf{Requirement} & \textbf{ASIL} & \textbf{Alloc.} & \textbf{Verification} \\ \midrule\endfirsthead
\toprule \textbf{ID} & \textbf{Requirement} & \textbf{ASIL} & \textbf{Alloc.} & \textbf{Verification} \\ \midrule\endhead
FSR-08.01 & While engaged, the item shall give the first driver alert within 6 s of the loss of attention, measured end to end from the driver's change of state to the alert (including camera frame time, DM model latency and policy timers), and shall reach the no-response reaction of FSR-02.09 within the time set in WP-S-04. & A & E-01/DM, E-01/HMI & T-SIL, T-VEH \\ \midrule
\bottomrule
\end{longtable}
}
\chapter{FTTI allocation}
Table~\ref{tab:ftti} splits each preliminary fault-tolerant time interval into a detection and a reaction budget for the technical safety requirements to meet. The behaviour today is kept in the record (WP-C-04 \S8).
{\small
\begin{longtable}{@{}p{3.6cm}p{2.5cm}p{2.9cm}p{2.9cm}p{2.4cm}@{}}
\caption{FTTI allocation (WP-C-04 \S8, preliminary).}\label{tab:ftti}\\
\toprule \textbf{Goal and case} & \textbf{FTTI} & \textbf{Detection} & \textbf{Reaction} & \textbf{Gap} \\ \midrule\endfirsthead
\toprule \textbf{Goal and case} & \textbf{FTTI} & \textbf{Detection} & \textbf{Reaction} & \textbf{Gap} \\ \midrule\endhead
SG-01 (continuous limiting) & $\le$ 0.5 s & 0 (per-frame check) & $\le$ 1 frame (10 ms) & — \\ \midrule
SG-01 (RX loss/corruption) & $\le$ 0.5 s & $\le$ 0.3 s & $\le$ 0.1 s & GAP-06 (partly closed), GAP-01 \\ \midrule
SG-01 (SoC command loss) & $\le$ 0.5 s & $\le$ 0.3 s & $\le$ 0.1 s & GAP-06 (partly closed) \\ \midrule
SG-01 (MCU execution fault) & $\le$ 0.5 s & $\le$ 0.1 s (HW watchdog; set by SG-07) & $\le$ 0.1 s & GAP-07 \\ \midrule
SG-02, SG-06 (loss without warning) & $\le$ 1 s to warning & $\le$ 0.5 s & $\le$ 0.5 s (warning) & GAP-16 \\ \midrule
SG-03, SG-04 & $\le$ 1 s & $\le$ 0.5 s & $\le$ 0.3 s & GAP-06 \\ \midrule
SG-05 (driver brake/cancel) & $\le$ 0.2 s & $\le$ 1 brake message period (25 ms at 40 Hz) & $\le$ 1 frame & GAP-46; verify on vehicle (T-VEH) \\ \midrule
SG-05 (steering override) & $\le$ 0.2 s & — & — & GAP-02 \\ \midrule
SG-07 (forwarding loss) & $\le$ 0.2 s (preliminary, WP-S-04 §3.6; working budget accepted, DQ-062) & $\le$ 0.1 s & $\le$ 0.1 s (relay release) & GAP-43, GAP-49 (\texttt{0xe7} stops forwarding); TSR-711 \\ \midrule
SG-07 (content) & per frame & 0 & — & GAP-42 (FSR-07.02) \\ \midrule
SG-08 (driver not supervising) & $\le$ 6 s from loss of attention to the first alert & $\le$ 5 s (DM policy, vision) & $\le$ 1 s (first alert) & GAP-21; time to the no-response reaction preliminary in WP-S-04 §3.8.4 \\ \midrule
\bottomrule
\end{longtable}
}
\chapter{Warning and degradation}
Figure~\ref{fig:ladder} shows the degradation levels from full function to lockout, what the driver is told at each, and the safe states they end in.
\begin{figure}[htbp]\centering
\includegraphics[width=\linewidth,height=0.8\textheight,keepaspectratio]{fig-c04-12-1-degradation-ladder.pdf}
\caption{The degradation ladder (WP-C-04 Figure 12-1).}\label{fig:ladder}
\end{figure}
\appendix
\chapter{Abbreviations}
Written from the project's single list of terms (\texttt{terms.py}), with only the terms this report uses.
\begin{description}\small
\item[ACC] adaptive cruise control
\item[AEB] automatic emergency braking
\item[AI] artificial intelligence
\item[AOU-nn] assumption of use (WP-C-01 §7); an R suffix marks the FSC's refined version
\item[ASIL] Automotive Safety Integrity Level, A (lowest) to D (highest), from ISO 26262
\item[DM] driver monitoring
\item[DQ-nn] decision of the project lead (decision register)
\item[E-nn] element of the item (WP-C-01 §3.1), E-01\ldots{}E-06
\item[EPS] electric power steering
\item[EXT-nn] external element credited only through an assumption of use (EXT-EPS, EXT-PCM, EXT-CLU)
\item[FCW] forward collision warning
\item[FI-nn] functional insufficiency (WP-C-06)
\item[FMEA] failure mode and effects analysis
\item[FMEDA] failure modes, effects and diagnostic analysis (hardware metrics)
\item[FSR-nn] functional safety requirement (FSC)
\item[FTA] fault tree analysis
\item[FTTI] fault-tolerant time interval: fault to hazard, the time a safety mechanism has
\item[GAP-nn] gap found by the gap assessment
\item[HMI] human-machine interface: display, sounds, cluster
\item[HW] hardware
\item[I0–I3] independence level of a review, I0 (none, e.g. the author or an AI) to I3 (independent organisation)
\item[ID] identifier
\item[ISO] International Organization for Standardization
\item[LKA] lane keeping assist
\item[MCU] microcontroller: the panda safety processor (STM32H7)
\item[ODD] operational design domain: where the system is meant to work
\item[OI-nn] open item of the document it appears in
\item[OPS-nn] operation step: the driver's pre-drive and field checks (WP-O-02 §3)
\item[PCM] powertrain control module (engine and drive)
\item[PCS] pre-collision system: Toyota's stock collision warning and braking
\item[QM] quality managed: no ASIL; ordinary quality processes suffice
\item[RX] receive
\item[SG-nn] safety goal (HARA §6)
\item[SoC] system on chip: the comma device's main processor, running openpilot (QM)
\item[SW] software
\item[T-HIL] verification by bench test with the real panda (hardware in the loop)
\item[T-SIL] verification by software test or process replay (WP-S-01 §1.1)
\item[T-VEH] verification by vehicle test
\item[TSR-nn] technical safety requirement (WP-S-02)
\item[TX] transmit
\item[WP-nn] work product, e.g. WP-C-04; the letter is its area (M management, C concept, S system, H hardware, W software, V validation, P processes, A analyses, O operation, K safety case)
\end{description}
\begin{thebibliography}{9}
\bibitem{iso26262-3} ISO 26262-3:2018, Road vehicles: Functional safety. Part 3: Concept phase. ISO, Geneva.
\end{thebibliography}
\end{document}
